\chapter*{Historical Background}

Before we can proceed to a direct analysis of the Gospel texts (after which point we will not go beyond them), there are still some observations we need to make regarding the real-life historical context behind the events. In the year 6 CE, after the death of King Herod the Great, Palestine lost the last remnants of its independence and was occupied by the Romans; two of the four historical provinces of Palestine, Judea and Samaria, fell under the rule of Roman procurators who were directly appointed by the imperial administration. This resulted in the sudden rise of a national liberation movement; its ideology was largely shaped by religious fundamentalists known as the Pharisees, but the best-organized force within the movement was a group of radical nationalists known as the Zealots:

\begin{smallquote}
    [The Zealots] agreed in all other things with the Pharisaicnotions; but they had an inviolable attachment to liberty... They also did not value dying any kinds of death; nor indeed did they heed the deaths of their relations and friends: nor could any such fear make them call any man lord.\footnote{\emph{Antiquities of the Jews}, Flavius Josephus, Book XVIII, Chapter 1, trans. William Whiston}
\end{smallquote}
It was the Zealots who played the decisive role in inciting the First Roman-Jewish
War of 66 CE, which ultimately brought the Jews to complete national catastrophe.

Traditionally, traveling preachers played a large role in shaping public opinion in Palestine. Historians are aware of about a dozen prophets who were preaching at around the same time as Jesus Christ and John the Baptist, and were comparable to them in popularity; almost all of them were executed by the Romans as potential leaders of a Jewish revolt. Nevertheless, the years leading up to the Jewish War saw 12 such revolts (not counting minor insurrections and acts of terrorism). The Romans reacted to these events in their usual manner: according to Flavius Josephus, they once crucified 2,000 people simultaneously in Jerusalem.

Balanced between these two irreconcilable forces was the local religio-political elite. While it initially grew out of a religious sect known as the Sadducees, it had completely dispensed with ideology by this point in time and was now a party of pragmatists. They were perfectly willing to collaborate with the Romans---or blood-sucking Martians, for that matter---to preserve the status quo; as soon as they got the chance, however, they would just as quickly paint themselves as the leaders of the national liberation movement (a trick the Communist partocracy of the former Soviet republics would later perform before our very eyes).

That actually happened, by the way. In Book II, Chapter 20 of \emph{The Jewish War}, Josephus describes the election of the new, ``revolutionary'' administration of the uprisen Jerusalem; they resulted in ``governorship of all affairs within the city'' being handed over to... the unsinkable high priest Ananus. Here is the commentary of Yakov Chertok, Russian translator and publisher of \emph{The Jewish War}, on this chapter:

\begin{smallquote}
    Accordingly, the results of the election proved quite unfavorable for the Zealots, despite the fact that they had gained decisive dominance over both the capital and the provinces following their defeat of Cestius and their expulsion of the Romans from the borders of the country. [...] Eleazar ben Simon, victor over Cestius and subsequently the main leader of the war, was completely passed over in the election; the even more powerful leader of the Zealots at that time, Eleazar ben Ananias, who gave the war its initial impetus [...], was given authority over the secondary province of Idumea, with the likely goal of sending him away from Jerusalem. \emph{The positions of greatest responsibility in Jerusalem and Galilee were granted to Roman allies who had previously gone into hiding to avoid the persecution of the Zealots} [emphasis mine---K.E.].
\end{smallquote}
Sounds familiar, don't you think?

These events came slightly later, though. In the time period we are concerned with, the Sadducee leadership found it more useful to demonstrate their loyalty to the ``imperial center.'' The unpopular regime, as is usually the case, flooded the country with secret police agents. Elsewhere in \emph{Antiquities of the Jews}, Josephus writes that from the time of Herod the Great (or rather, the Bloody), the torture, execution, and ``disappearance'' of dissidents became common practice:

\begin{smallquote}
    And many there were who were brought to the citadel Hyrcania, both openly, and secretly; and were there put to death. And there were spies set everywhere, both in the city, and in the roads: who watched those that met together.\footnote{Book XV, Chapter 10.}
\end{smallquote}

In response, a militant faction of Zealots known as the Sicarii launched a campaign of terror against the occupiers and their local collaborators, the crowning achievement of which was their assassination of the high priest Jonathan. The Zealots had extensive clandestine networks at their disposal, and their agents permeated every part of the state apparatus. In addition, they were in frequent contact with bands of robbers (or ``partisan detachments,'' if you will), which the country was quite literally swarming with. Some of these armed groups numbered up to several hundred men; the legendary field commander Eleazar, for example, terrorized the outskirts of Jerusalem for almost 20 years, outlasting several procurators and high priests.

Likewise, there is no doubt that the Parthian Empire's intelligence services were also actively operating in Syria and Palestine; during this time, the Parthians were putting up a tough and rather effective resistance to the Roman\emph{Drang nach Osten}.\footnote{ In his book \emph{The Craft of Intelligence}, longtime CIA director Allen Dulles highlights the role that Parthian intelligence played in this ``strategy of containment''; thanks in large part to the personal efforts of King Mithridates VI, it became one of the best intelligence agencies in the entire history of the ancient world. (K.E.)} To be perfectly clear: I have yet to see any concrete evidence of foreign support of the Jewish national liberation movement. That being said, it is hard to imagine that such sophisticated politicians as the Parthians would have been unaware of the basic tactical principle, ``the enemy of my enemy is my friend.''

In short, the real Palestine was, to put it in modern-day spy lingo, ``a country with a complex agent-operational situation,'' kind of like Lebanon or El Salvador in the 1980s. It lived up quite poorly to the idyllic picture painted of it by the Gospels: a land of shady olive groves where the people spend most of their time having edifying conversations and debating religion and philosophy.

I felt the need to go on this historical excursion for one reason, and one reason only. Whenever anybody (McDowell included), referring to Christ's Judea, casually writes that ``the authorities did this'' or ``the authorities wanted that,'' that is pure nonsense. There were really \emph{two} sets of authorities (as well as a clandestine, albeit quite powerful one), and their interests, while sometimes overlapping on certain issues, were completely different on the whole. To top it all off, the general instability of the situation clearly hindered these authorities' ability to maintain internal unity. In these conditions, the natural rivalry between groups, or the individuals within them, could take the form of open conflict, combined with the most unexpected and unnatural of temporary alliances (``Who are we going to be friends against today?'').\footnote{ A common Russian phrase. (Z.B.)} In this regard, let us recall the state of affairs of the political leadership and intelligence services in Nazi Germany, which was so vividly reconstructed in the dearly beloved Soviet television series \emph{Seventeen Moments of Spring}.

Anyone who thinks this is purely conjecture should find Chertok's commentary on Chapter 13 of \emph{The Jewish War} quite interesting:


\begin{smallquote}
    The high priest Jonathan facilitated the appointment of Felix as procurator, as a result of which he incurred the wrath of the Sicarii. Meanwhile, however, Felix began to grow weary of Jonathan, who frequently reproached him for his cruel and unjust acts, and he wished to be rid of him. To that end, he came to an agreement with the Sicarii, who, despite being enemies of Felix, nevertheless offered him their services to assassinate the high priest whom they equally despised.
\end{smallquote}

Biblical scholars have long highlighted a strange fact: even though Jesus spoke out just as sharply against the decadent Sadducee elite as he did against the dogmatist Pharisees, not once did he say a single word about the Zealots, whose activities were one of the most pressing issues of the day. Taking into account that the Zealots held the greatest power and authority in Jesus' native Galilee, the English biblical scholar S. G. F. Brandon argued in his monographs \emph{Jesus and the Zealots} and \emph{The Trial of Jesus of Nazareth} that Jesus, if not an outright member of the Zealot movement, was at the very least quite sympathetic to it.

In any case, at least one of the Apostles is reliably known to have been a Zealot: Simon (Lk 6:15). However, in his book \emph{The State in the New Testament}, the German biblical scholar Oscar Cullmann provides evidence that four other Apostles were members of the Zealot movement: Peter, Judas, James the Great, and John. A telling detail: during the Last Supper, when Christ makes an allegorical plea for his Apostles to ``sell their garments and buy [a sword],'' they take him literally and reply, ``Lord, look, here are two swords'' (Lk 22:38).\footnote{Scripture taken from the New King James Version®. Copyright © 1982 by Thomas Nelson. Used by permission. All rights reserved.} For what it's worth, Judea wasn't exactly Texas, and it was completely unheard of to carry weapons out in the open.

In terms of the problem before us, there is absolutely no need to dive into speculation as to what organizational ties Jesus might have had to the national liberation movement. What is important for us here is another aspect entirely. In times of social upheaval, any socially significant figure (a class to which Jesus certainly belonged---as did the other prophets of his time), regardless of what their own plans and desires may be, becomes a \emph{political} figure. And that means they become either an independent actor, or an object of manipulation by outside forces.